\section{Описание}

\subsection{Источник данных}
Корпус построен из статей русской Википедии (\texttt{ru.wikipedia.org}).
Статьи получаются через открытый \MYhref{https://www.mediawiki.org/wiki/API:Main_page}{MediaWiki Action API}:
метод \texttt{list=random} (с фильтрами: основное пространство имён, без
перенаправлений) возвращает случайные статьи, метод \texttt{action=parse} ---
HTML-код статьи. Ключ API и авторизация не нужны; запросы отправляются с
идентифицирующим \texttt{User-Agent}, как требуют правила Викимедиа.

\subsection{Характеристики корпуса}
Документ --- «сырая» HTML-страница статьи (вёрстка MediaWiki: заголовки
\texttt{h2}/\texttt{h3}, абзацы \texttt{p}, списки, таблицы, инфобоксы,
навигационные блоки, ссылки на источники) плюс JSON-файл с метаданными: url,
источник, идентификатор и заголовок статьи, очищенный текст, число символов и
слов, время загрузки (Unix time stamp). Язык --- русский (доля кириллицы среди
букв $\approx$94{,}2\%, остальное --- латиница в названиях и терминах).
Тематика разнообразная: населённые пункты, персоналии, биология, техника,
спорт, история и т.д.

\subsection{Разбиение на документы}
Принцип «одна статья --- один документ». Файлы хранятся как
\texttt{corpus/wikipedia/<pageid>.html} и \texttt{<pageid>.json}, где
\texttt{pageid} --- идентификатор статьи в Википедии. Повторный запуск
программы пропускает уже скачанные статьи. Статьи короче 1000 символов
(заготовки) в корпус не включаются.

\subsection{Выделение текста}
Из HTML текст выделяется библиотекой BeautifulSoup. Сначала удаляются
служебные элементы: \texttt{style}, таблицы, сноски (\texttt{sup.reference}),
списки источников, навигационные блоки (\texttt{.navbox}), пометки
\texttt{.hatnote}, миниатюры изображений и ссылки <<править>>. Затем
собирается текст абзацев, заголовков и элементов списков.

\subsection{Существующие поисковики}
Для поиска по выбранным документам уже можно использовать:
\begin{itemize}
    \item встроенный поиск Википедии
    (\MYhref{https://ru.wikipedia.org/wiki/Служебная:Поиск}{Служебная:Поиск}
    и API \texttt{list=search});
    \item Google и Яндекс с ограничением по сайту:
    \texttt{site:ru.wikipedia.org}.
\end{itemize}

\subsection{Примеры запросов и недостатки выдачи}
Запросы выполнены через встроенный поиск Википедии (API \texttt{list=search}).
\begin{enumerate}
    \item \texttt{компьютер мышь} --- найдено 1864 статьи; первой идёт
    <<Компьютерная мышь>>, но следом --- страница значений <<Мышь (значения)>> и
    общие статьи (<<Периферийное устройство>>, <<Персональный компьютер>>), которые
    лишь упоминают оба слова.
    \item \texttt{<<информационный поиск>>} (точная фраза) --- найдено 65
    статей; первой идёт <<Информационный поиск>>, дальше --- статьи по смежным
    темам (кластеризация документов, метапоиск). Ранжирование по релевантности
    не объясняется, пользователь не видит, почему статья оказалась выше.
    \item \texttt{яблоко} --- найдено 8857 статей; в начале выдачи плоды,
    политическая партия <<Яблоко>>, сорт <<Золотое яблоко>> и анатомический
    термин <<Адамово яблоко>>: многозначность слова не разрешается, нужно
    уточнять запрос вручную.
\end{enumerate}

Общие недостатки: нет учёта морфологии и синонимов в явном виде, не разрешается
многозначность, нет возможности настраивать ранжирование, а внешние поисковики
(Google, Яндекс) возвращают страницы вне нашего корпуса.

\subsection{Статистика по корпусу}
Статистика получена командой \texttt{python3 main.py -{}-stats} по полностью
скачанному корпусу.

\begin{tabular}{|l|c|}
\hline
Показатель & Значение \\ \hline
Размер \enquote{сырых} данных & 2074 МБ ($\approx$2{,}0 ГБ) \\ \hline
Количество документов & 30027 \\ \hline
Размер выделенного текста & 231{,}5 МБ \\ \hline
Средний размер документа & 69{,}1 КБ \\ \hline
Средний объём текста в документе & 7{,}7 КБ \\ \hline
Среднее число слов в документе & 639 \\ \hline
\end{tabular}

\pagebreak
